\chapter*{คำนำ}
\addcontentsline{toc}{chapter}{คำนำ}

จุลชีววิทยาการเกษตร (Agricultural Microbiology) เป็นศาสตร์ที่มีความสำคัญอย่างยิ่งยวดต่อความมั่นคงทางอาหาร สิ่งแวดล้อม และระบบนิเวศเกษตรที่ยั่งยืน ในยุคปัจจุบันที่ภาคการเกษตรทั่วโลกกำลังเผชิญหน้ากับความท้าทายหลายประการ ทั้งภาวะการเสื่อมโทรมของดิน การแพร่ระบาดของโรคและแมลงศัตรูพืช การเปลี่ยนแปลงสภาพภูมิอากาศ (Climate Change) และความต้องการลด ละ เลิก การใช้สารเคมีสังเคราะห์ทางการเกษตร เพื่อก้าวเข้าสู่โมเดลเศรษฐกิจชีวภาพ เศรษฐกิจหมุนเวียน และเศรษฐกิจสีเขียว (Bio-Circular-Green Economy: BCG Model) จุลินทรีย์จึงเป็นคำตอบทางธรรมชาติและกุญแจดอกสำคัญในการฟื้นฟูความอุดมสมบูรณ์ของดิน เพิ่มผลผลิตพืช และสร้างสมดุลเชิงนิเวศ

ตำราเล่มนี้เรียบเรียงขึ้นโดยมีวัตถุประสงค์เพื่อใช้เป็นเอกสารหลักในการจัดการเรียนการสอนรายวิชาจุลชีววิทยาการเกษตร สำหรับนักศึกษาระดับปริญญาตรีและบัณฑิตศึกษา สาขาวิชาเกษตรศาสตร์ จุลชีววิทยา วิทยาศาสตร์ชีวภาพ และเทคโนโลยีการเกษตร ตลอดจนเป็นแหล่งอ้างอิงทางวิชาการสำหรับคณาจารย์ นักวิจัย นักวิชาการเกษตร และเกษตรกรก้าวหน้าที่มุ่งเน้นการทำเกษตรแม่นยำและเกษตรยั่งยืน

โครงสร้างเนื้อหาของตำราเล่มนี้ได้รับการออกแบบอย่างเป็นระบบ แบ่งออกเป็น 9 บทเรียนหลัก และ 1 บทนำ ครอบคลุมตั้งแต่หลักการพื้นฐานจนถึงเทคโนโลยีระดับแนวหน้า
\begin{itemize}
    \item \textbf{บทนำ} ประวัติวิวัฒนาการ หมุดหมายสำคัญทางประวัติศาสตร์ และกรอบแนวคิดการประยุกต์ใช้จุลินทรีย์ในการขับเคลื่อนเศรษฐกิจชีวภาพ
    \item \textbf{บทที่ 1} ความหลากหลาย โครงสร้างประชากร และแหล่งอาศัยระดับไมโครของจุลินทรีย์ในดิน
    \item \textbf{บทที่ 2} วัฏจักรชีวธรณีเคมีและการหมุนเวียนธาตุอาหารหลัก (ไนโตรเจน คาร์บอน ฟอสฟอรัส และกำมะถัน)
    \item \textbf{บทที่ 3} ปฏิสัมพันธ์เชิงบวกระหว่างจุลินทรีย์กับรากพืช (Rhizosphere, PGPR และ Mycorrhiza)
    \item \textbf{บทที่ 4} การควบคุมศัตรูพืชและโรคพืชโดยชีววิธี (Biocontrol Agents) ด้วยไตรโคเดอร์มา บาซิลลัส และเชื้อราก่อโรคแมลง
    \item \textbf{บทที่ 5} สาเหตุ พยาธิสภาพ และการวินิจฉัยโรคพืชที่เกิดจากจุลินทรีย์พร้อมแนวทางการจัดการแบบบูรณาการ
    \item \textbf{บทที่ 6} จุลินทรีย์ในการจัดการของเสียทางการเกษตร การทำปุ๋ยหมักชีวภาพ และการผลิตพลังงานก๊าซชีวภาพ
    \item \textbf{บทที่ 7} เทคโนโลยีการหมักในอุตสาหกรรมเกษตรแปรรูป โพรไบโอติกสำหรับสัตว์ และผลิตภัณฑ์ชีวภาพมูลค่าสูง
    \item \textbf{บทที่ 8} บทบาทของจุลินทรีย์ต่อการกักเก็บคาร์บอน การบรรเทาผลกระทบจากภาวะโลกร้อน และการพัฒนาปุ๋ยชีวภาพ
    \item \textbf{บทที่ 9} การประยุกต์ใช้เทคโนโลยีดีเอ็นเอรหัสแท่ง (Metagenomics) ชีวสารสนเทศศาสตร์ และปัญญาประดิษฐ์ (AI) ในการประเมินดัชนีสุขภาพดิน
\end{itemize}

ผู้เขียนหวังเป็นอย่างยิ่งว่า ตำราเล่มนี้จะเป็นประโยชน์ในการเสริมสร้างองค์ความรู้ ปูพื้นฐานการวิจัย และจุดประกายความคิดสร้างสรรค์ทางวิชาการ เพื่อนำเทคโนโลยีจุลินทรีย์ไปประยุกต์ใช้ยกระดับภาคการเกษตรของประเทศให้เจริญก้าวหน้า มั่นคง และยั่งยืนสืบไป

\vspace{1.5cm}
\begin{flushright}
    \textbf{ผู้ช่วยศาสตราจารย์ ดร.จิรภัทร จันทมาลี}\\
    สาขาวิชาจุลชีววิทยา คณะวิทยาศาสตร์และเทคโนโลยี\\
    มหาวิทยาลัยราชภัฏรำไพพรรณี\\
    กันยายน 2569
\end{flushright}
\clearpage
